\section{The latest reduction of the plus height to two character estimates}
\label{sec:sr-two-counts}
Anton's subsequent
\href{https://github.com/antoshashakov/riemann-hypothesis/blob/76ffabc62ab67b126afc128013ab6efa6f7d2a7c/verification/TwoCharacterCounts.lean}{\srcid{TwoCharacterCounts}}
reduces the polynomial plus-height input to two rational-character
bounds. We give the precise estimates and their assembly here.
Neither estimate is asserted unconditionally in this section.
We first supply the trivial-character vanishing needed in that assembly.

\subsection{Two vanishing arguments over \(\mathbb Q\)}
\begin{srproposition}\label{prop:sr-rational-regular}
For \(\ell=3,7\) and every \(a\geq1\),
\(T_{\ell,a}(\mathbb Q)=0\).
\end{srproposition}
\begin{proof}
Restriction takes this strict group into
\(T_{\ell,a}(\mathbb Q(\zeta_\ell))\), which is zero by
Proposition~\ref{prop:sr-regularsha}.
The extension has degree \(\ell-1\), a unit in \(\mathbb Z_\ell\).
Restriction is injective: restriction followed by corestriction is
multiplication by \(\ell-1\), or equivalently the finite-quotient
cohomology kernel vanishes by averaging over that quotient.
The same argument applies to the local restriction conditions.
Thus the original group is zero.
\end{proof}

\begin{srproposition}[The resonant trivial character]\label{prop:sr-rational-resonant}
For every odd prime \(\ell\) and every positive \(a\) with
\(\ell-1\mid2a\), one has \(T_{\ell,a}(\mathbb Q)=0\).
\end{srproposition}
\begin{proof}
Let \(\omega:\operatorname{Gal}(\mathbb Q(\mu_{\ell^\infty})/\mathbb Q)
\to\mu_{\ell-1}\subset\mathbb Z_\ell^\times\) be the finite
Teichmuller part of the cyclotomic character.
The condition implies that \(\omega^a\) has order at most two.
Its fixed field \(K\) in \(\mathbb Q(\zeta_\ell)\) is either
\(\mathbb Q\) or the unique quadratic subfield.
In the quadratic case its discriminant is
\((-1)^{(\ell-1)/2}\ell\).
Indeed the Gauss-sum calculation in Section~\ref{sec:sr-dyadic-dual}
works for every odd prime: it gives
\(\tau_\ell^2=(-1)^{(\ell-1)/2}\ell\), and this squarefree radicand
is \(1\bmod4\).
The integral quadratic basis \(1,(1+\sqrt D)/2\) has discriminant
\(D\); the trace and norm integrality conditions give the usual
full ring of integers for squarefree \(D\equiv1\bmod4\).

In both cases \(\ell\nmid h_K\). For the quadratic case we include
an elementary bound. Minkowski gives an integral ideal in each class
of norm at most \(B\leq\sqrt\ell\).
There are at most \(d(n)\) ideals of norm \(n\): at a split rational
prime the count is the exponent plus one, and the inert and ramified
counts are smaller. Here \(d(n)\) is the divisor function.
For \(B\geq1\),
\[
 h_K\leq\sum_{n\leq B}d(n)
 =\sum_{u\leq B}\lfloor B/u\rfloor
 \leq B(1+\log B)
 \leq\sqrt\ell(1+\log\sqrt\ell)<\ell .
\]
The last inequality is \(\log z<z-1\) for \(z>1\).
For \(K=\mathbb Q\) the class number is one.

Let \(K_\infty\) be the subfield of
\(\mathbb Q(\mu_{\ell^\infty})\) fixed by
\(\ker(\omega^a)\) in its finite factor \(\mu_{\ell-1}\).
Then \(K_\infty/K\) is a \(\mathbb Z_\ell\)-extension, and
\(\chi_{\rm cyc}^{-a}\) is trivial on \(G_{K_\infty}\).
There is a unique prime above \(\ell\) in \(K\), totally ramified
through this tower, and no other finite prime ramifies.
The one-ramified-prime class-number lemma also applies to these
cyclic degree-\(\ell\) steps if \(K\) is real.
For odd \(\ell\) no real place can complexify in such a Galois
extension. In its proof the unit representation is then
\(\mathbb Q[G]^{r_1(K)+r_2(K)}-\mathbb Q\);
the same Herbrand quotient \(1/\ell\) and ideal calculation apply.
Thus \(\ell\) is prime to every finite-layer class number.

The strict descent of Proposition~\ref{prop:sr-regularsha} now applies
to this tower, as can also be seen directly.
A strict cocycle has finite image. In a sufficiently high layer the
coefficient action fixes that image and the chosen local primitive
at the unique prime above \(\ell\).
Its restriction is a homomorphism, zero at that decomposition group
and unramified everywhere else.
Artin reciprocity and the prime-to-\(\ell\) class number kill it.
The cocycle therefore descends to
\(\operatorname{Gal}(K_\infty/K)\simeq\mathbb Z_\ell\).
The generator acts on \(\mathbb Q_\ell/\mathbb Z_\ell\) by a scalar
different from one, and multiplication by the difference is
surjective. Subtracting a coboundary kills its value at that
generator, hence kills the cocycle by continuity.
We obtain \(T_{\ell,a}(K)=0\).
Finally \([K:\mathbb Q]\leq2\) is prime to \(\ell\), so the
restriction argument of Proposition~\ref{prop:sr-rational-regular}
gives the result over \(\mathbb Q\).
\end{proof}

\subsection{The rational groups and their normalizations}
For an odd prime \(\ell\), choose a finite splitting extension of
\(\mathbb Q_\ell\), with ring \(\mathcal O\), fraction field \(Q\),
and \(r_{\mathcal O}=v_{\mathcal O}(\ell)>0\).
For a Dirichlet character \(\psi\) of level \(N\), let
\(\mathcal V'_{\ell,N,a}(\psi)\) be the subgroup of
\[
 H^1(\mathbb Q,(Q/\mathcal O)(\psi^{-1}\chi_{\rm cyc}^{-a}))
\]
whose restrictions to all finite inertia groups and to the full
local groups at primes dividing \(\ell N\) are zero.
These are restrictions of cohomology classes, not assertions that
every representative cocycle is pointwise zero.
Put
\[
 v_{\ell,a}(\psi)=
 \frac{\operatorname{length}_{\mathcal O}
              \mathcal V'_{\ell,N,a}(\psi)}{r_{\mathcal O}}
                  \quad\hbox{in }[0,\infty].
\]
This is the normalized length called \srcid{Vn} when its finiteness
has been proved; \srcid{VLe} includes that finiteness as well as
the weighted inequality. An infinite length is never set to zero.

For a principal host let \(H\subseteq(\mathbb Z/N)^\times\)
be its fixing subgroup, and let
\[
 \mathcal X_a=\{\psi:\psi|_H=1,\ \psi(-1)=(-1)^a\}.
\]
All these characters have \(2\)-power order and
\(|\mathcal X_a|\leq[E:\mathbb Q]\leq N\).
The plus parity here is different from the minus parity in
Proposition~\ref{prop:sr-PB}.

The restriction and averaging proof in that proposition, applied
to the other parity, gives
\begin{equation}\label{eq:sr-plus-character}
 s^+_{\ell,E}(a)\leq
             \sum_{\psi\in\mathcal X_a}v_{\ell,a}(\psi).
\end{equation}
There is no index loss: every odd \(\ell\) is prime to
\([E:\mathbb Q]\).
We check the extra local conditions used in \(\mathcal V'\).
At \(p\ne\ell\), if \(\psi\) is ramified, a nontrivial inertia
value differs from one by an \(\mathcal O\)-unit, because its
order is prime to \(\ell\).
The argument in Section~\ref{sec:minus-characterwise} turns an
inertia coboundary into a full local coboundary.
If \(\psi\) is unramified, an unramified local cocycle factors
through the procyclic unramified quotient.
Its Frobenius scalar is \(\psi(\operatorname{Frob}_p)^{-1}p^{-a}\ne1\):
a positive rational power \(p^a>1\) cannot be a root of unity.
Subtract a coboundary using surjectivity of this nonzero scalar
minus one on \(Q/\mathcal O\). The cocycle is then zero at
Frobenius, hence zero by continuity.
Thus every relevant unramified class is already locally zero
at these additional primes. This proves that imposing the
\(N\)-prime conditions introduces no missing local term in
\eqref{eq:sr-plus-character}.

For \(\psi=1\), the rational vanishing propositions imply
\[
 v_{\ell,a}(1)=0\quad
 \hbox{if }\ell\in\{3,7\}\ \hbox{or }\ell-1\mid2a .
\]
The statements remain true over \(\mathcal O\), since its scalar
extension is finite free over \(\mathbb Z_\ell\), and continuous
cohomology and the local kernels commute with finite direct sums.
Adding the further local conditions only takes a subgroup.

\subsection{The two remaining estimates}
Write \(M=Nt\), with \(t\geq2\).
All the following quantifiers concern the same sufficiently large
principal hosts as \eqref{eq:sr-principal-height}.
The fixed thresholds and the exponents precede the host and \(t\).
The splitting coefficient ring may vary with \(\ell\).

\begin{srdefinition}\label{def:sr-character-inputs}
The generic estimate asks for one natural exponent \(A\) such that,
for each \(1\leq a<t\) and every finite prime set \(S\) with
\(\ell\notin\{2,3,7\}\) and \(\ell-1\nmid2a\),
\begin{equation}\label{eq:sr-generic-count}
 \sum_{\ell\in S}(\log\ell)
       \sum_{\psi\in\mathcal X_a}v_{\ell,a}(\psi)\leq M^A .
\end{equation}
The nontrivial estimate asks for one natural exponent \(B\) such that
\begin{equation}\label{eq:sr-nontrivial-count}
 (\log\ell)v_{\ell,a}(\psi)\leq(\ell N a)^B
\end{equation}
whenever \(\psi\in\mathcal X_a,\ \psi\ne1,\ \ell\) is odd, and
\[
 \ell\in\{3,7\}\quad\hbox{or}\quad
                       (a\geq2\ \hbox{and }\ell-1\mid2a).
\]
Finiteness is included in both estimates.
\end{srdefinition}
The first is a simultaneous sum over arbitrarily large finite
prime sets, rather than a pointwise bound which could not be
summed over all primes. These two estimates are still assumptions
of the next proposition.

\begin{srproposition}\label{prop:sr-two-counts}
Definition~\ref{def:sr-character-inputs} implies the polynomial
plus-height bound \eqref{eq:sr-principal-height}.
One may use the uniform exponent \(A+3B+5\).
At the exceptional pair \(3,7\) alone, exponent \(2B+2\) suffices.
\end{srproposition}
\begin{proof}
All sums below may first be taken over an arbitrary fixed finite
prime set. The inputs and the vanishing results make each rational
length finite; \eqref{eq:sr-plus-character} then proves finiteness
of each plus length.

For \(\ell\notin\{2,3,7\}\) and a fixed \(a\), the generic primes
contribute at most \(M^A\).
Every other prime satisfies \(\ell-1\mid2a\), so
\(\ell\leq2a+1\leq2t\), and necessarily \(a\geq2\).
There are at most \(2t\) such primes.
The trivial character contributes zero.
For every other character,
\[
 \ell N a\leq2Nt^2=2Mt\leq M^3,
\]
since \(t\leq M\) and \(M\geq2\).
Its weighted contribution is therefore at most \(M^{3B}\).
There are at most \(N\) characters.
Summing the fixed-\(a\) estimate over \(a<t\), and using \(tN=M\),
bounds the entire outside contribution by
\[
 t(M^A+2tN M^{3B})
 \leq M^{A+1}+2M^{3B+2}
 \leq3M^{A+3B+2}\leq M^{A+3B+4}.
\]
All exponents are nonnegative; the last inequality uses \(M^2\geq4\).

At \(\ell=3,7\), the trivial character again contributes zero for
every positive \(a\).
Here \(M\geq7\) for the principal hosts, so
\(\ell N a\leq M^2\).
Each remaining weighted contribution is at most \(M^{2B}\).
The two primes, at most \(N\) characters and fewer than \(t\)
twists give at most
\[
          2tN M^{2B}=2M^{2B+1}\leq M^{2B+2}.
\]
Add the inside and outside estimates.
Since \(A+3B+4\geq2B+2\) and \(M\geq2\), their sum is at most
\(M^{A+3B+5}\).
Equation~\eqref{eq:sr-plus-character} proves the asserted plus bound.
The constants are independent of the chosen finite prime set,
host, splitting rings and twist range, as required.
\end{proof}

In the source's notation the exceptional set at conductor one is
\(B_\chi(1)=\{3,7\}\).
Its quadratic auxiliary field is \(\mathbb Q(\sqrt{-7})\), whose
Minkowski bound \((2/\pi)\sqrt7<2\) gives class number one.
The first eligible quadratic auxiliary prime is seven and the
other auxiliary prime is three; the base character has order one.
The odd prime factors in the defining exceptional set therefore
reduce to those of \(7\cdot3\), together with three itself.
An odd-order nontrivial host character cannot occur, because its
order divides a power of two.
This is the elementary content of the exceptional-set and
Case III simplifications in the contributors' files.
The proof above avoids retaining the case labels: the generic,
resonant and exceptional conditions are stated directly.

The arithmetic content not supplied here is
\eqref{eq:sr-generic-count} and~\eqref{eq:sr-nontrivial-count}.
Their source proofs involve the global counting and geometric
comparisons in Part II. The vanishing and the assembly above
are not substitutes for those two estimates.

